\section{Lõi RISC-V gọn nhẹ và tích hợp bộ tăng tốc}
\subsection{PicoRV32 và giao tiếp bus gốc}
PicoRV32 là lõi RISC-V hướng tới diện tích nhỏ, có các lựa chọn giao tiếp bộ nhớ gốc, AXI4-Lite và Wishbone. Lõi thực hiện lệnh qua nhiều chu kỳ, không phải lõi đơn chu kỳ; thời gian còn phụ thuộc cấu hình và đáp ứng của bộ nhớ. Giao tiếp gốc dùng cơ chế valid/ready, phù hợp để kết nối RAM và ngoại vi qua bộ giải mã địa chỉ~\cite[README: Variants; PicoRV32 Native Memory Interface; Cycles per Instruction]{YosysHQ2024}.

Trong MyTinyMobileNetSoC, lựa chọn có ý nghĩa là giữ CPU làm chủ giao dịch và đưa các phép tính ra ngoại vi MMIO. Firmware dùng lệnh tải/lưu thông thường để nạp thanh ghi, khởi phát tính toán và đọc kết quả; không sử dụng giao diện lệnh đồng xử lý PCPI cho các toán tử CNN. Kiến trúc này cho phép quan sát rõ chi phí phục vụ dữ liệu, nhưng mỗi toán hạng phải đi qua CPU. RAM dùng chung chứa mã và tensor không đồng nghĩa ngoại vi được quyền tự đọc RAM, như sẽ mô tả ở Chương~\ref{chap:kientruc}.

Tài liệu PicoRV32 công bố tài nguyên và tần số cho những cấu hình lõi cụ thể~\cite[README: Features; Performance]{YosysHQ2024}. Những kết quả đó không đại diện cho SoC có RAM, bộ giải mã và các datapath CNN. Do đó, luận văn đánh giá tài nguyên và timing của toàn bộ thiết kế sau định tuyến ở Chương~\ref{chap:toiuu}, thay vì áp thông số lõi độc lập cho hệ thống tích hợp.

\subsection{Rocket/RoCC như một phương án đối chiếu}
Trong hệ sinh thái Rocket Chip, RoCC cho phép bộ tăng tốc nhận lệnh tùy chỉnh cùng toán hạng thanh ghi và trả kết quả về CPU. Bộ tăng tốc có thể dùng giao diện tới cache L1 hoặc kết nối TileLink tới hệ thống nhớ; lựa chọn đường truy cập ảnh hưởng thông lượng và độ phức tạp tích hợp~\cite{chipyardRoCC}. Cần phân biệt đường \emph{điều khiển} qua lệnh RoCC với đường \emph{dữ liệu} qua cache hoặc cổng nhớ: có RoCC không tự động có nghĩa mọi thiết kế đều dùng cùng một cơ chế DMA.

Khác với phương án đó, MyTinyMobileNetSoC không mở rộng tập lệnh để điều khiển CNN và không có DMA ở các ngoại vi. Cách tiếp cận này đổi lại sự đơn giản của giao diện phần mềm bằng chi phí tải/lưu và điều phối trên PicoRV32. Ngược lại, khả năng chuyển khối dữ liệu của bộ đồng xử lý đòi hỏi thêm điều khiển bộ nhớ và bộ đệm. Vì vậy, lựa chọn CPU trong chương được xem xét cùng phương thức tích hợp; danh mục pipeline hay hiệu năng lõi riêng không đủ để lựa chọn kiến trúc tăng tốc.
